\documentclass[a4paper]{article}
\usepackage[pdftex]{graphicx}
\usepackage[utf8]{inputenc}
\usepackage{enumerate}
\usepackage{siunitx}
\sisetup{locale=DE} 
\usepackage{href-ul}
\hypersetup{
	colorlinks=true,
	linkcolor=blue,
	urlcolor=blue}
\usepackage{icomma}
\usepackage{amssymb}
\usepackage{geometry}
\geometry{a4paper, top=15mm, left=15mm, right=15mm, bottom=15mm,
	headsep=10mm, footskip=12mm}

\begin{document}

\begin{enumerate}[1.]

{\bf \item Was versteht man unter $\dots$}

\begin{enumerate}[a)]
\item $\dots$ gedämpfter Schwingung
\item $\dots$ Amplitude einer Schwingung
\end{enumerate}

%Aufgabe 2
{\bf \item Ein Pendel führt in 5 Sekunden 20 Schwingungen durch. Ermitteln Sie die Frequenz des Pendels.}

%Aufgabe 3 
{\bf \item Ein Federpendel (Federhärte $D=\SI{1,9}{\newton\over\meter}$, Masse $m=\SI{200}{\gram}$) befindet sich zum Zeitpunkt $t=0$ im unteren Umkehrpunkt.}

\begin{enumerate}[a)]
\item Berechnen Sie die Schwingungsdauer dieses Federpendels auf Zehntelsekunden genau. Geben Sie auch die Frequenz an.
\item Erläutern Sie anhand einer Skizze, wo sich der schwingende Pendelkörper nach 
$t_1=\SI{3,5}{\second}$ befindet und in welche Richtung er sich bewegt.
\end{enumerate}

%aufgabe 4 
{\bf \item Bei einem Federpendel ( Pendelmasse $m$, Federhärte $D$ ) soll die Schwingungsdauer halbiert werden.} Nennen Sie zwei verschiedene Möglichkeiten.


%Aufgabe 5
{\bf \item Die Schwingung eines Federpendels ( Federhärte $\SI{0,49}{\newton\over\meter}$, Amplitude $\SI{12}{\centi\meter}$)}
startet zum Zeitpunkt $t=0$ im unteren Umkehrpunkt und erreicht zwei Sekunden später den oberen Umkehrpunkt. 

\begin{enumerate}[a)]
\item Wie groß ist die Schwingungsdauer?
\item Wie groß ist die angehängte Masse?
\item Wie groß ist die Rückstellkraft und die Beschleunigung im oberen Umkehrpunkt ?
\item An welchen ersten drei Zeitpunkten ist die Geschwindigkeit des Pendelkörpers am größten?
\end{enumerate}

{\bf \item Bei einer anderen harmonischen Schwingung} mit der Amplitude $ \hat{x}=\SI{10}{\centi\meter} $ misst man eine Schwingungsdauer von  $\SI{0,4}{\second}$ Erläutern Sie, welche Frequenz diese Schwingung bei einer Amplitude von $ \hat{x}=\SI{20}{\centi\meter} $ besitzt.
%Aufgabe 3
{\bf \item Ermitteln Sie aus den unten stehenden Diagrammen Amplitude, Wellenlänge, Schwingungsdauer} und die Frequenz sowie die Ausbreitungsgeschwindigkeit der mechanischen Welle.

\includegraphics[width=12 cm]{sinus95.pdf}

\end{enumerate} 

\fbox{
	\begin{minipage}{0.5\textwidth}
		Zur Lösung bitte \href{https://www.okuyakl.de/phys/p10pebL408/ll408.pdf}{hier klicken} oder den QR-Code scannen.\\
	Weitere Arbeitsblätter gibt es unter 
	
	\href{https://www.okuyakl.de}{www.okuyakl.de}
	\end{minipage}
	\hfill
	\begin{minipage}{0.4\textwidth}
		\includegraphics[width=1.5 cm]{../../viecher/zwe03}
		\includegraphics[width=3 cm]{qr408}
		\includegraphics[width=2 cm]{../../viecher/afanticon1}
		
	\end{minipage}}

\end{document}
